\documentclass{article}

\usepackage{iftex}
\ifPDFTeX
  \usepackage[T1]{fontenc}
  \usepackage[utf8]{inputenc}
\fi
\usepackage[UTF8]{ctex}
\usepackage[backend=biber,style=gb7714-2015,sorting=gb7714-2015,gbalign=left,sorting=gb7714-2015]{biblatex}
\usepackage{amsmath, amsthm, amssymb, amsfonts}
\usepackage{thmtools}
\usepackage{graphicx}
\usepackage{setspace}
\usepackage{geometry}
\usepackage{float}
\usepackage{hyperref}
\usepackage[english]{babel}
\usepackage{framed}
\usepackage[dvipsnames]{xcolor}
\usepackage{tcolorbox}
\usepackage{tikz-cd}
\usepackage{tikz} 

\tcbuselibrary{breakable}

\addbibresource{references.bib}

\colorlet{LightGray}{White!90!Periwinkle}
\colorlet{LightOrange}{Orange!15}
\colorlet{LightGreen}{Green!15}

\newcommand{\HRule}[1]{\rule{\linewidth}{#1}}

\declaretheoremstyle[name=Theorem,]{thmsty}
\declaretheorem[style=thmsty,numberwithin=section]{theorem}
\tcolorboxenvironment{theorem}{colback=White}

\declaretheoremstyle[name=Definition,]{prosty}
\declaretheorem[style=prosty,numberlike=theorem]{definition}
\tcolorboxenvironment{definition}{colback=White}

\declaretheoremstyle[name=Lemma,]{prcpsty}
\declaretheorem[style=prcpsty,numberlike=theorem]{lemma}
\tcolorboxenvironment{lemma}{colback=White}

\newtheorem{proposition}[theorem]{Proposition}
\newtheorem{corollary}[theorem]{Corollary}
\newtheorem{remark}[theorem]{Remark}
\newtheorem{example}[theorem]{Example}
\newtheorem{exercise}[theorem]{Exercise}


\def\E{\mathop{\mathcal E}}
\def\D{\mathop{\mathcal D}}
\def\lhdeq{\mathop{\trianglelefteq}}
\def\<{\mathop{\langle}}
\def\>{\mathop{\rangle}}
\def\l{\mathop{\ell}}
\newcommand{\g}{\mathfrak{g}}  
\newcommand{\h}{\mathfrak{h}} 
\newcommand{\n}{\mathfrak{n}}  
\newcommand{\U}{\mathcal{U}}  
\newcommand{\C}{\mathbb{C}} 
\newcommand{\borel}{\mathfrak{b}}
\newcommand{\Ind}{\mathrm{Ind}} 
\newcommand{\Hom}{\mathrm{Hom}}  
\newcommand{\Ext}{\mathrm{Ext}}

\newcommand{\slC}{\mathfrak{sl}(2, \mathbb{C})}
\newcommand{\vir}{\mathfrak{vir}}
\newcommand{\der}{\mathrm{d}}

\newcommand{\Z}{\mathbb{Z}}
\newcommand{\End}{\mathrm{End}}
\newcommand{\N}{\mathbb{N}}

% ------------------------------------------------------------------------------

\begin{document}

% ------------------------------------------------------------------------------
% Cover Page and ToC
% ----------------------------------------------------------------------

\section*{About ideal structure of Weyl algebra and introduction to differential operators.}

\begin{center}
Xinyuan. Zhang
\end{center}

Let $ K[x_1, \ldots , x_n]: = K[X] $. Denote $ \hat{x_i}f=x_i \cdot f $ for $ f \in K[x_1, \ldots , x_n] $. The Weyl algebra $ A_n $ is defined as
\[
A_n= < \hat{x_1}, \ldots, \hat{x_n}, \partial_1, \ldots, \partial_n  >
\]
The relation is $ [\partial_i, x_j] = \delta_{ij} $.

This algebra has a $ K $-basis 
\[
x^{\alpha} \partial^{\beta}=
\] 

We begin with the ideal of $ A_n $.

For any $ D \in A_n $. $ \deg (D)=|\alpha|+|\beta| $. 

\begin{theorem} 
$ \forall D,D' \in A_n $, we have:
\begin{enumerate}
    \item $ \deg(D+D') \le \mathrm{max}\{ \deg(D), \deg(D') \} $.
    \item $ \deg(D \cdot D') = \deg(D) + \deg(D') $.
    \item $ \deg ([DD']) \le \deg (D) + \deg (D')-2 $. 
\end{enumerate}
\end{theorem}

\begin{proof} 
1
\end{proof}

\begin{theorem} 
$A_n$ is simple (as associative algebra).
\end{theorem}

\begin{proof} 
2
\end{proof}


\begin{corollary}
All non-zero elements in $ \operatorname{End}(A_n)  $ are injective. 
\end{corollary}

For characteristic $ p > 0 $. 

\section{Rings of differential operators}

Let $ R $ be a commutative $ k $-algebra. For any $ a\ in K $ Define map $ v \mapsto av $. order 0-n.

\begin{lemma}
Let $ D'=\operatorname{Der}(R)+R $ 
\end{lemma}

\begin{proof} 
3
\end{proof}

\begin{definition} 
$ D(R) = \bigcup_{} $ 
\end{definition}

\begin{proposition}
If $ P \in D^n $ and $ Q \in D^m $, then $ PQ \in D^{m+n} $.   
\end{proposition}

\begin{proof} 
4
\end{proof}

\begin{proposition}
$ \operatorname{Der}(K[X]) $ 
\end{proposition}

\begin{proof} 
5
\end{proof}

We want to prove $ A_n=D(K[X]) $.

\begin{lemma}
$ P \in D(K[X]) $, $ \forall i, [P,x_i]=0 $, then $ P \in K[X] $.   
\end{lemma}

\begin{proof} 
6
\end{proof}

\begin{definition} 
$C_r$
\end{definition}

\begin{lemma}
For $ P_1, \ldots, P_n \in C_{r-1} $, $ \forall [P_i,x_j]=[P_j, x_i] $  
\end{lemma}

\begin{proof} 
7
\end{proof}

\begin{theorem} 
$ D(K[X])=A_n $ 
\end{theorem}

\begin{proof} 
8
\end{proof}

\section{$ A_n $-modules}

\begin{definition} 
torsion elements
\end{definition}

\begin{definition} 
torsion module
\end{definition}

\begin{lemma}
If $ M $ is an irreducible $ R $-module.
\begin{enumerate}
    \item $ M \cong R/\operatorname{Ann}(u) $.
    \item If $ R $ is not a division ring, then $ M $ is a torsion module.
\end{enumerate}
\end{lemma}

\begin{proposition}
$ K[X] $ is an irreducible torsion $ A_n $-module. Further more, $ K[X] \cong A_n / \sum_{i=1}^n A_n \partial_i $. 
\end{proposition}

\begin{proof} 

\end{proof}

We also have $ A_n/ \sum A_n x_i \cong K[\partial_1, \ldots, \partial_n]$.

\begin{proposition}
Twisting $ R $-mod $ M $.  
\end{proposition}

\begin{proposition}
(1) M irr $ \Leftrightarrow $ M twis irr 
(2) tor $ \Leftrightarrow $ twis tor
(3) $ J $ is a left ideal of $ R $ so is $ \sigma(J) $.   
\end{proposition}

Let $ F: A_n \to A_n $, $ x_i \mapsto \partial_i $, $ \partial_i \mapsto -x_i $. We call $ M_F $ the $ F $-transform of $ M $.

\begin{proposition}
$ (K[X])_F=K[\partial_1, \ldots, \partial_n]=K[\partial] $ 
\end{proposition}




































































\end{document}
